\section{向量: 可运算的 ``点''} \label{sec:vector}

我们用数量来描述一些具体事物的特征, 比如长度, 温度, 价格等.
数量可以进行加减乘除等运算, 以描述事物的联系与变化.
而弹幕制作中涉及的许多事物, 比如子弹的位置, 速度等,
需要用多个数值来表示.
为了用运算描述它们, 我们引入\emph{向量}的概念.

\subsection{向量是什么}

初学者可能会对``向量是什么''感到困惑,
这不奇怪, 向量的概念在数学史上经历了相当丰富的演化过程,
限于篇幅原因我无法介绍向量的发展过程.
为了明确向量是什么, 我们从弹幕制作的角度给向量下一个明确的定义.

% 一开始, 向量的思想隐含在几何学和力学的研究中,
% 主要用来表示力, 速度, 加速度等物理量.
% 这时向量被视为``带有方向的线段'',
% 可以说是一个具体的几何对象.

% 后来向量的概念逐渐从几何中解放出来, 向代数范畴靠近.
% 这时向量主要被视为由若干数值组成的一个量,
% 通过代数运算规则建立了一套向量分析体系,
% 主要服务于电磁学研究.

% 再后来, 数学家不再满足于用几个数值构建向量,
% 而是将其抽象为具有特定加法和数乘运算规则的任何元素,
% 以至于很多和几何上的向量完全不搭边的东西也可以视为向量,
% 比如函数.
% 这个对我们弹幕制作就有点过于抽象了, 唉抽代是这样的.

\begin{definition}[向量]
  设有$n$个数值$a_1,a_2,...,a_n$,
  将数组$(a_1,a_2,...,a_n)$称为一个\emph{$n$维向量}.
  $n$称为该向量的\emph{维度},
  $a_i$ ($i=1,2,...,n$)称为该向量的第$i$个\emph{分量}.
\end{definition}

我们一般用英文字母和希腊字母来指代一个数量, 比如$x,y,\theta$等.
一般我们会用上方加箭头的英文字母或希腊字母指代一个向量,
比如$\vec r=(3,4)$表示一个二维向量$\vec r$, 它的值为$(3,4)$.

在几何上, 我们主要用向量来描述两种几何对象: 点和有向线段.
如图\ref{fig:vector-geometry},
对于平面直角坐标系上的一点$P$,
我们将它的直角坐标作为它对应的向量的值,
并且仍然用$P$表示该向量 (因为这并不会造成混淆).
对于由点$A$指向点$B$的有向线段,
我们用$\vv{AB}$表示它对应的向量,
该向量的值为从$A$移动到$B$的平移量的直角坐标.

注: 将有向线段整体进行平移不会改变它对应向量的值.
如果有需要, 比如描述一条直线激光, 可以考虑用起点位置和平移量
两个向量来进行描述.

\begin{figure}[htbp]
  \centering
  \begin{tikzpicture}
    \draw[-latex] (-1.5,0) -- (3,0)
    node[right] {$x$};
    \draw[-latex] (0,-0.5) -- (0,2)
    node[left] {$y$};

    \node[below left] at (0,0) {$O$};
    \draw[->] (0.5,0.5) node[left] {$A$}
    -- (2.5,1.5) node[right] {$B$}
    node[below right,pos=0.5] {$\vv{AB}=(2,1)$};
    \node[dot,label={left:$P=(-0.5,0.5)$}] at (-0.5,0.5) {};
  \end{tikzpicture}
  \caption{向量的几何应用}
  \label{fig:vector-geometry}
\end{figure}

除直角坐标以外, 我们还使用极坐标来描述位置.
为了方便表述, 我们也可以用极坐标表示向量,
但这时向量的值是对应的直角坐标而不是极坐标本身.
我的意思是, 当我们说一个向量的值为$\braket{r,\theta}$,
它是一种简写, 这个向量实际的值是$(r\cos\theta,r\sin\theta)$.

\subsection{向量的几何属性}

还是要强调一下, 本教程所讲的向量,
纯粹是几个数值所构成的一个代数对象,
它的存在不依赖于有向线段等几何对象.
我们所说向量的几何属性,
只是由构成向量的几个数值进行一些特定的计算,
只不过这些特定计算在几何中有着明确的应用而已.

\subsubsection{长度和方向}

把向量$\vec{a}$对应到$n$维直角坐标系中的有向线段,
该有向线段的长度称为向量$\vec{a}$的\emph{长度},
也称为向量的\emph{大小}或\emph{模长}, 记作$|\vec{a}|$.

设$\vec{a}=(a_1,...,a_n)$,
该向量的长度可以用
\[
  \vec{a} = \sqrt{a_1^2 + \cdots + a_n^2}
\]
计算. 这是勾股定理在$n$维空间的推广.

向量$\vec{a}$对应的有向线段的方向称为该向量的\emph{方向}.
如果$\vec{a}$是二维向量$(a_1,a_2)$,
那么可以用对应有向线段的\emph{方向角}来表示$\vec{a}$的方向,
将该方向角记作$\angle(\vec{a})$, 它的一个等价角可以用下式计算:
\[
\angle(\vec{a}) = \Angle(a_1,a_2).
\]

长度为1的向量称为\emph{单位向量}.

长度为0的向量, 或者说, 各分量都为0的向量, 称为\emph{零向量},
我们用$\vec0$或者直接用0表示.
长度不为0的向量则称为\emph{非零向量}.

\begin{definition}(二维向量的夹角)
  对于两个非零的二维向量$\vec r_1, \vec r_2$,
  设$\vec r_1$对应一个有向线段$\vv{OA}$,
  $\vec r_2$对应一个有向线段$\vv{OB}$.
  将始边为$OA$, 终边为$OB$的任意角
  称为$\vec r_1,\vec r_2$的\emph{夹角},
  记作$\angle(\vec r_1,\vec r_2)$.
  并且规定夹角的取值范围为
  $-180\degree < \cdot \le 180\degree$.

  注: 本教程采用的夹角与通常的夹角概念稍有不同,
  通常的夹角概念是不管始边和终边的顺序的, 也不会出现负数夹角.
\end{definition}

\subsubsection{相等和相反}

对两个相同维度的向量
$\vec{a}=(a_1,...,a_n),\vec{b}=(b_1,...,b_n)$,
如果两向量的各分量对应相等, 即
\[a_1=b_1,a_2=b_2,...,a_n=b_n,\]
我们说这两个向量是\emph{相等}的,
记作$\vec{a}=\vec{b}$.

如果各分量对应互为相反数, 即
\[a_1=-b_1,a_2=-b_2,...,a_n=-b_n,\]
那么我们说这两个向量互为\emph{相反向量},
记作$\vec{a}=-\vec{b}$.

从几何角度考虑, 对于两个相同维度的向量$\vec{a},\vec{b}$,
\begin{itemize}
  \item 若$\vec{a},\vec{b}$的大小相等, 方向相同, 那么$\vec{a}=\vec{b}$;
  \item 若$\vec{a},\vec{b}$的大小相等, 方向相反, 那么$\vec{a}=-\vec{b}$.
\end{itemize}

注: 对于任意两个点$A,B$, 自然地,
$A=\vv{OA}, \vv{AB}=-\vv{BA}$总是成立.

\subsubsection{平行(共线)与垂直}

设有两个非零向量$\vec{a},\vec{b}$,
如果它们的方向相同或相反, 那么称$\vec{a},\vec{b}$
\emph{平行}, 又称\emph{共线}, 记作
$\vec{a}\parallel\vec{b}$;
如果它们的方向互相垂直, 那么称$\vec{a},\vec{b}$
\emph{垂直}, 记作
$\vec{a}\bot\vec{b}$.
规定零向量与任意向量平行且垂直.

\begin{remark}
  向量的平行和共线完全没有区别, 从有向线段的角度来看, 这可能显得有些奇怪.
  如图\ref{fig:parallel}, $AB$与$EF$共线, $CD$与$EF$平行且长度相等.
  由于向量不关注有向线段的起点和终点, 只关注长度和方向,
  从向量的角度, $\vv{CD}$与$\vv{EF}$没有任何区别
  (它们的横纵坐标分量是对应相等的), 所以没有办法区分平行和共线.
  在一些场景中需要区分平行和共线, 这时就需要把线段的起点位置也考虑进来.
  
  \begin{figure}[htbp]
    \centering
    \begin{tikzpicture}
      \draw[->] (0,0) node[left] {$A$}
      -- (15:2) node[right] {$B$};
      \draw[->] (15:4) node[left] {$E$}
      -- +(15:1) node[right] {$F$};
      \draw[->] (3,-0.5) node[left] {$C$}
      -- +(15:1) node[right] {$D$};
    \end{tikzpicture}
    \caption{平行与共线}
    \label{fig:parallel}
  \end{figure}
\end{remark}
